\section{Problem Setup}
\label{sec:setup}

% MANUSCRIPT-NODE: SETUP-P01
We analyze the actor-side dynamics generated by repeatedly reusing signed feedback from a historical update distribution. Let $\nu$ denote a distribution over state--action pairs, representing an offline dataset, a replay buffer, or trajectories collected by earlier or stale policies, and let $\widehat A(s,a)$ be the advantage-like coefficient supplied to the actor. During one actor step, both $\nu$ and $\widehat A$ are treated as fixed with respect to $\theta$. The resulting empirical actor field and discrete update are
\begin{equation}
\begin{aligned}
\mathbf F(\theta)
&=\mathbb E_{(s,a)\sim\nu}\!\left[\widehat A(s,a)\nabla_\theta\log\pi_\theta(a\mid s)\right],\\
\theta_{t+1}&=\theta_t+\eta\mathbf F(\theta_t).
\end{aligned}
\label{eq:signed-field}
\end{equation}
This field is not assumed to equal the exact on-policy policy gradient, because $\nu$ need not match the current state--action distribution~\cite{sutton1999policy,levine2020offline}. Define the positive and negative components
\begin{equation}
\widehat A^+=\max(\widehat A,0),\qquad \widehat A^-=\max(-\widehat A,0),
\label{eq:sign-split}
\end{equation}
so that $\mathbf F$ decomposes into attraction from $\widehat A^+$ and repulsion from $\widehat A^-$. For one negative sample, the update magnitude factorizes as
\begin{equation}
\left\|\widehat A^-\nabla_\theta\log\pi_\theta(a\mid s)\right\|_2
=\widehat A^-\left\|\nabla_\theta\log\pi_\theta(a\mid s)\right\|_2.
\label{eq:influence-factorization}
\end{equation}
The first factor measures disadvantage severity; the second is learner-relative score geometry. At the aggregate level, sample frequency and directional coherence further determine whether these contributions cancel or reinforce. Freezing $\nu$ and $\widehat A$ is therefore a mechanism-isolation device: it asks what the actor score alone can do under historical reuse, without claiming that critic error, importance ratios, or jointly evolving actor--critic dynamics are absent in full algorithms.
% END-MANUSCRIPT-NODE: SETUP-P01

% MANUSCRIPT-NODE: SETUP-P02
To compare continuous and categorical policies without forcing them to share the same amplification law, we use a common learner-relative coordinate. For a fixed context $s$ and historical action $a$, define the negative log probability
\begin{equation}
D_\theta(s,a)=-\log\pi_\theta(a\mid s).
\label{eq:remoteness}
\end{equation}
For a Gaussian policy, $D_\theta$ equals a scale-dependent constant plus one half of the squared Mahalanobis distance from the policy mean. For a categorical policy, it is exactly the selected action's surprisal. In sequence models the same quantity may be evaluated at token level or normalized across a completion, but the normalization must be fixed before comparing examples. Crucially, remoteness is dynamic: the historical action and its label remain fixed while $D_\theta(s,a)$ changes as the learner moves. We separately define the squared score response
\begin{equation}
R_\theta(s,a)=\left\|\nabla_\theta\log\pi_\theta(a\mid s)\right\|_2^2.
\label{eq:score-response}
\end{equation}
The notation is shared, but no universal relation between $D_\theta$ and $R_\theta$ is assumed. Gaussian mean coordinates can have unbounded response, categorical direct-logit scores are bounded, and shared neural parameters introduce an additional Jacobian. Proposition~\ref{prop:score-remoteness} therefore derives the policy-family-specific static relation before we analyze repeated reuse.
% END-MANUSCRIPT-NODE: SETUP-P02
